\chapter{Inference rules}

\section{Notations and Terminology}

\paragraph{Judgements and terms} To describe the type system, its types and
the inference rules, this document uses "judgements" on "terms".\\*
For example, $\tau:T$ is a judgement telling that $\tau$ is a term of type $T$.

\paragraph{Assumptions} Zero or more assumptions (comma separated) can be
used for a judgement; they're specified before the "turnstile" symbol
$\then$.\\*
For example, $\EnvSym \then \tau:T$ can be read as "if $\EnvSym$ is true, then
$\tau$ is a term of type $T$".

\paragraph{Premise} Judgements can be made based on zero or more premises
(semicolon separated judgements).\\*
For example, $\infrule{\EnvSym_1\then\tau_1:T\semiq\EnvSym_2\then\tau_2:T}{
\EnvSym_1,\EnvSym_2\then \code{$\tau_1$ + $\tau_2$}: T}$ means that with a
premise that both judgements $\EnvSym_1\then\tau_1:T$ and
$\EnvSym_2\then\tau_2:T$ are true, then
$\EnvSym_1,\EnvSym_2\then\code{$\tau_1$+$\tau_2$}:T$ is true.

\paragraph{Code} Text written with \code{monospaced} font denotes the text
as written in code verbatim. Some types may also be written with
\texttt{Monospace} font (usually in PascalCase).

If the code evaluates to some value, to denote its type the
"$:\texttt{SomeType}$" notation will be used.\\*
For example: $\code{59 + 10} : \Int$

To denote to what the value evaluates to, the $\eval$ symbol will be used.\\*
For example: $\code{59 + 10} \eval 69$

\paragraph{Functions} Functions are denoted by a list of zero or more input
arguments followed by an output argument.\\*
For example, $(\tau_1,\tau_2)\to\tau_3$ describes a function type that takes
two arguments of type $\tau_1$ and $\tau_2$ respectively, and produces a value
of type $\tau_3$.

\clearpage

\paragraph{"Generic" types} Some types are "generic" (a.k.a. existential
types), meaning that they declare one or more type variable (a.k.a. type
parameters).\\*
The notation $T[\tau:=\sigma]$ describes a generic type $T$ with a type
parameter $\tau$ of some type $\sigma$.\\*
The notation $t : \exists\sigma.T$ means that $t$ is a term of concrete
type $T$ with the type parameter $\tau$ equal to $\sigma$.
\todo{tipi esistenziali? Rivedere notazione}

\paragraph{List} To denote a list type, the existential type notation will
be used $\List[\nu=\Type]$.\\*
To denote a value, for example of type $\exists\Int.\List$, this notation will
be used: $\glist{1,2,3}$

\paragraph{Map} To denote a map type, the existential type notation will
be used $\Map[\kappa:=\Type,\nu:=\Type]$.\\*
To denote a value, for example of type $\exists\Text.\exists\Int.\Map$, this
notation will be used: $\map{a\mapsto1,b\mapsto2,c\mapsto3}$

\paragraph{Identifiers set} A set $\Id=\set{i\mid i:=\text{identifier}}$ is
used to encompass all possible identifiers allowed by the grammar.
\todo{formalizzare}

\paragraph{Product types}

\subparagraph{Tuple} A tuple type is denoted with the cartesian product
operator.\\*
For example, $\tau_1 \times \tau_2 \times \tau_3$ describes a 3-tuple type
with three fields of type $\tau_1, \tau_2, \tau_3$.

\subparagraph{Record} The notation $\tuple{x_1:\tau_1, x_2:\tau_2}$ describes a
record type (a tuple whose fields are named).

\subparagraph{Struct} The notation $(\xi, \tuple{x_1:\tau_1, x_2:\tau_2})$
describes a struct type (a pair of an identifier $\xi$ and a record type).
\todo{rivedere notazione struct}

\begin{definition}[Types widening and narrowing]
	A type $\eta$ is said to be "narrower" than a type $\omega$ iff a value
	described by $\eta$ can also be described by $\omega$ without any loss of
	information\footnotemark[1].\\*
	Conversely, a type $\omega$ is said to be "wider" than a type $\eta$ iff it
	can describe any value described by $\eta$ without loss of information.\\*
	A "narrower than"\footnotemark[2] relationship is denoted using
	the $\narreq$ operator. For example: $\eta\narreq\omega$.
\end{definition}

\footnotetext[1]{
	In this language, a loss of information can be a loss of data, for example
	when coercing a \Real~to an \Int, but also a loss of "semantic data".
	See~\See{sec:product_types_compatibility}.
}

\footnotetext[2]{
	As you'll see in \See{par:compatibility_reflexivity}, the symbol $\narreq$
	actually denotes a "narrower or equal than" relationship; that's why it is
	used instead of the more common symbol $<:$ which is reserved to denote a
	"strictly narrower than" relationship.
}

\paragraph{Least Wider Bound} The $\lwb$ operator produces the least wide type
for the operand types
\[
	\infrule[aligned]{
		\EnvSym & \then \tau_1:\Type \semiq \tau_2:\Type \\*
		\EnvSym & \then \exists_{\omega:\Type}{\tau_1\narreq\omega \land \tau_2\narreq\omega} \\*
		\EnvSym & \then \nexists_{\omega'\Type}{\tau_1\narreq\omega' \land \tau_2\narreq\omega' \land \omega'\narreq\omega}
	}{
		\EnvSym\then \tau_1 \lwb \tau_2 = \omega
	}
\]
\todo{Least Upper Bound?}

\clearpage

\subsection{Type environments}

The definition of an expression type is done within a "context" (a.k.a. type
environment) which describes the type of any identifier or any other
expressions appearing in the described expression.

The general type environment is denoted by $\EnvSym$, a function that maps an
identifier to its type.

To avoid ambiguities and to more clearly describe expressions, $\EnvSym$ is
partioned in the $\EnvVar,\EnvConst,\EnvFunc,\EnvStucts$ environments:

\begin{itemize}
	\item $\EnvVar$: the Variable environment, a function that maps a variable
		identifier $i$ to its type $\tau$
		\[
			\begin{split}
				\EnvVar : \Id &\to \Type\\*
				\EnvVar(i)    &\mapsto \tau
			\end{split}
		\]
	\item $\EnvConst$: the Constants environment, a function that maps a
		constant identifier $i$ to its type $\tau$
		\[
			\begin{split}
				\EnvConst : \Id &\to \Type\\*
				\EnvConst(i)    &\mapsto \tau
			\end{split}
		\]
	\item $\EnvFunc$: the Functions environment, a function that maps a
		function identifier $i$ to its definition signature $(i, \phi)$
		\[
			\begin{split}
				\EnvFunc : \Id &\to (\Id, \Function)\\*
				\EnvFunc(i)    &\mapsto (i, \phi)
			\end{split}
		\]
		With a type $\tau$, it maps to the function defined within the type
		$\tau$
		\[
			\begin{split}
				\EnvFunc : (\Type,\Id) &\to (\Id, \Function)\\*
				\EnvFunc(\tau, i)      &\mapsto (i, \phi)
			\end{split}
		\]
	\item $\EnvStucts$: the Structs environment
		\[
			\begin{split}
				\EnvStucts : \Id &\to (\Id, \Record)\\*
				\EnvStucts(i)    &\mapsto (i, \rho)
			\end{split}
		\]
		A function that maps a typename $i$ (struct definition identifier) to
		its type definition (identifier $i$ and record $\rho$)
	% \item $\EnvModules$: the Modules environment
	% 	\[
	% 		\begin{split}
	% 			\EnvModules : \Id & \to \Module\\*
	% 			\EnvModules(i)    & \mapsto m
	% 		\end{split}
	% 	\]
	% 	A function that maps an identifier $i$ to the module
\end{itemize}

\todo{formalizzare structs, funzioni e moduli}

The $\EnvSym$ function can now be defined as follows
\[
	\begin{split}
		\EnvSym : \Id & \to \Type\\*
		\EnvSym(i) & \mapsto \begin{cases}
			\EnvStucts(i) &\text{ if } \EnvVar(i)=\varnothing \land \EnvConst(i)=\varnothing \land \EnvFunc(i)=\varnothing \\*
			\EnvFunc(i) &\text{ if } \EnvVar(i)=\varnothing \land \EnvConst(i)=\varnothing \\*
			\EnvConst(i) &\text{ if } \EnvVar(i)=\varnothing \\*
			\EnvVar(i) & \\*
			% \EnvModules(i) & \text{ if } \\*
		\end{cases}
	\end{split}
\]
which formally describes the lookup order and priority of a symbol $i$.

\paragraph{Environment extension} Given how scopes and symbols lookup work, it
will be useful to define "new" type environments extending an existing
environment. The notation $E\env{i:\tau}$ will be used as a shorthand to denote
$E \cup \env{i:\tau}$, which is a new environment that inherits from $E$ and
extends it with $\env{i:\tau}$. Formally:
\[
	E\env{i:\tau}(j) \mapsto \begin{cases}
		i &\text{ if } j = i \\*
		E(j) & \\*
	\end{cases}
\]

\clearpage

\section{Types}

\nb{
	The syntax for expressing a type or an expression may be ambiguous without
	context. For this reason the \code{type.} prefix syntax is used, but note
	that in some contexts (e.g. variable declaration) the \code{type.} is not
	needed.
}

\todo{
	TODO: menzionare uso di type system strutturale e nominale

	TODO: usare 0-based indexes

	TODO: usare notazione tipi polimorfici invece che solo Type
}

\subsection{Primitive types}

\[
	\begin{split}
		&\EnvSym\then \Void:\Type \\
		&\EnvSym\then \Bool:\Type \\
		&\EnvSym\then \Natural:\Type \\
		&\EnvSym\then \Int:\Type \\
		&\EnvSym\then \Real:\Type \\
		&\EnvSym\then \Text:\Type \\
	\end{split}
\]

\subsection{Collection types}

\[
	\begin{split}
		&\EnvSym\then \List[\nu=\Type]:\Type \\
		&\EnvSym\then \Map[\kappa:=\Type,\nu:=\Type]:\Type \\
	\end{split}
\]

\[
	\infrule{
		\EnvSym\then \tau:\Type
	}{
		\EnvSym\then \code{type.[$\tau$]} \eval \exists\tau.\List
	}
\]

\[
	\infrule{
		\EnvSym\then \tau:\Type \semiq \sigma:\Type
	}{
		\EnvSym\then \code{type.[$\tau$ => $\sigma$]} \eval \exists\tau.\exists\sigma.\Map
	}
\]

\nb{
	When the type parameters are not needed in the rule, they will be omitted.
	For example: $\tau=\List$, $\expr:\Map$
}

\clearpage

\subsection{Product types}

\paragraph{Tuple}

\[
	\EnvSym\then \Tuple : \Type
\]

\nb{
	When the whole definition of a tuple type is not needed, the type will
	simply be denoted by \Tuple.
}

\[
	\infrule{
		\EnvSym\then n \geq 0 \semiq \forall_{i\in[1,n]}{\tau_i:\Type}
	}{
		\EnvSym\then \code{tuple <$\tau_1$, $\dotsc$, $\tau_n$>} \eval \tau_1 \times \dotsb \times \tau_n
	}
\]

\paragraph{Record}

\[
	\EnvSym\then \Record : \Type
\]

\nb{
	When the whole definition of a record type is not needed, the type will
	simply be denoted by \Record.
}

\[
	\infrule{
		\EnvSym\then n \geq 0 \semiq \forall_{i\in[1,n]}{\tau_i:\Type, x_i \in \Id}
	}{
		\EnvSym\then \code{record <$x_1$ $\tau_1$, $\dotsc$, $x_n$ $\tau_n$>} \eval \tuple{x_1:\tau_1, \dotsc, x_n:\tau_n}
	}
\]

\paragraph{Struct}

\[
	\EnvSym\then \Struct : \Type
\]

\nb{
	When the whole definition of a struct type is not needed, the type will
	simply be denoted by \Struct.
}

\[
	\infrule{
		\EnvSym\then \xi \in \Id \semiq n \geq 0 \semiq \forall_{i\in[1,n]}{\tau_i:\Type, x_i \in \Id}
	}{
		\EnvSym,\code{struct $\xi$ <$x_1$ $\tau_1$, $\dotsc$, $x_n$ $\tau_n$>} \then \code{$\xi$} \eval (\xi,\tuple{x_1:\tau_1,\dotsc,x_n:\tau_n})
	}
\]

\todo{che notazione usare per \Struct???}

\clearpage

\subsection{Functions}

\[
	\EnvSym\then \Function : \Type
\]

\nb{
	When the whole definition of a function type is not needed, the type will
	simply be denoted by \Function.
}

\[
	\infrule{
		\EnvSym\then n \geq 0 \semiq \forall_{i\in[1,n]}{\tau_i:\Type,x_i\in\Id} \semiq \tau_{n+1}\Type
	}{
		\EnvSym\then \code{::($x_1$ $\tau_1$, $\dotsc$, $x_n$ $\tau_n$) => $\tau_{n+1}$} \eval (\tau_1,\dotsc,\tau_n)\to\tau_{n+1}
	}
\]

\[
	\infrule{
		\EnvSym\then n \geq 0 \semiq \forall_{i\in[1,n]}{\tau_i:\Type,x_i\in\Id}
	}{
		\EnvSym\then \code{::($x_1$ $\tau_1$, $\dotsc$, $x_n$ $\tau_n$)} \eval (\tau_1,\dotsc,\tau_n)\to\Void
	}
\]

\[
	\infrule{
		\EnvSym\then \xi \in \Id \semiq n \geq 0 \semiq \forall_{i\in[1,n]}{\tau_i:\Type,x_i\in\Id} \semiq \tau_{n+1}\Type
	}{
		\EnvSym,\code{$\xi$($x_1$ $\tau_1$, $\dotsc$, $x_n$ $\tau_n$) => $\tau_{n+1}$ \{\}} \then \code{$\xi$} : (\tau_1,\dotsc,\tau_n)\to\tau_{n+1}
	}
\]

\[
	\infrule{
		\EnvSym\then \xi \in \Id \semiq n \geq 0 \semiq \forall_{i\in[1,n]}{\tau_i:\Type,x_i\in\Id}
	}{
		\EnvSym,\code{$\xi$($x_1$ $\tau_1$, $\dotsc$, $x_n$ $\tau_n$) \{\}} \then \code{$\xi$} : (\tau_1,\dotsc,\tau_n)\to\Void
	}
\]

% \subsection{Range}

\todo{Range. Tipo concreto non implementato ancora}

\subsection{Compatibility rules}

\todo{menzionare covarianze (tipi composti e prodotto) e
controvarianza (parametri funzioni?)}

\paragraph{Reflexivity}
\label{par:compatibility_reflexivity}
\[
	\infrule{
		\EnvSym\then \tau:\Type
	}{
		\EnvSym\then \tau \narreq \tau
	}
\]

\paragraph{Transitivity}
\[
	\infrule{
		\EnvSym\then \tau_1 \narreq \tau_2 \semiq \tau_2 \narreq \tau_3
	}{
		\EnvSym\then \tau_1 \narreq \tau_3
	}
\]

\subsubsection{Primitive types}

\[
	\infrule{
		\EnvSym\then \tau : \Type
	}{
		\EnvSym\then \Void \narreq \tau
	}
\]

\[
	\EnvSym\then \Natural \narreq \Int
	\qquad \EnvSym\then \Int \narreq \Real
\]

\subsubsection{Collection types}

\[
	\infrule{
		\EnvSym\then \tau_1 = \List[\nu_1:=\Type]\\*
		\EnvSym\then \tau_2 = \List[\nu_2:=\Type]\\*
		\EnvSym\then \nu_1 \narreq \nu_2
	}{
		\EnvSym\then \tau_1 \narreq \tau_2
	}
\]

\[
	\infrule{
		\EnvSym\then \tau_1 = \Map[\kappa_1:=\Type, \nu_1:=\Type]\\*
		\EnvSym\then \tau_2 = \Map[\kappa_2:=\Type, \nu_2:=\Type]\\*
		\EnvSym\then \kappa_1 \narreq \kappa_2 \land \nu_1 \narreq \nu_2
	}{
		\EnvSym\then \tau_1 \narreq \tau_2
	}
\]

\clearpage

\subsubsection{Product types}
\label{sec:product_types_compatibility}

Since the widening and narrowing is crucial for type inference, Tuples, Records
and Structs require some additional attention as they have some differences.

\vspace{\baselineskip}

In order to quickly understand why the widening between product types behaves
the way it does, ask yourself: \textit{what information does this data
carry?}\\*
For example, \code{<1,2,3,"foo">} is a tuple with three integer fields and a
text field; no additional semantic information is given, and those values could
mean anything.\\*
Instead, \code{<x:1, y:2, z:3, name:"foo">} is a record with three integer
fields and a text field, with the additional information about their name;
their semantic meaning is still uncertain.\\*
Finally, \code{Rgb.<x:1, y:2, z:3, name:"foo">} is a struct describing a color
which has an integer field for each red, green and blue channel, and a field
representing the common name of the color.\\*
Instead, \code{Point.<x:1, y:2, z:3, name:"foo">} is a struct describing
a point in a coordinate system which has a three integer field for each
dimension and a label.

\vspace{\baselineskip}

While all previous examples evaluates to an "equal" value (the data they
carry), their semantic meaning is potentially very different. This means that,
while the tuple or the record could be widened to either \code{Rgb} or
\code{Point} (where the "added information" is their semantic meaning), a
\code{Point} cannot be losslessly converted to \code{Rgb} or other narrower
types, because we would "lose" their semantic meaning.\\*
If you're ok with losing such information, conversion between "incompatible"
product types can still be achieved via explicit type coercion (casting).

\clearpage

\paragraph{Tuple}

\[
	\infrule{
		\EnvSym\then n \geq 0 \semiq \forall_{i\in[1,n]}{\tau_i:\Type} \semiq \theta_1 = \bigtimes_{i \in [1,n]}{\tau_i}\\*
		\EnvSym\then m \geq n \semiq \forall_{i\in[1,m]}{\omega_i:\Type} \semiq \theta_2 = \bigtimes_{j \in [1,m]}{\omega_j}\\*
		\EnvSym\then \forall_{i \in [1,n]}{\tau_i \narreq \omega_i}
	}{
		\EnvSym\then \theta_1 \narreq \theta_2
	}
\]
A tuple type $\theta_1$ is narrower than the tuple type $\theta_2$ if the first
$n$ fields of $\theta_1$ have a type $\tau$ narrower than the first $n$ fields
types $\omega$ of $\theta_2$.

\vspace{\baselineskip}

\[
	\infrule{
		\EnvSym\then n \geq 0 \semiq \forall_{i\in[1,n]}{\tau_i:\Type} \semiq \theta = \bigtimes_{i \in [1,n]}{\tau_i}\\*
		\EnvSym\then m \geq n \semiq \forall_{j\in[1,m]}{x_j \in \Id, \omega_j:\Type} \semiq \rho = \tuple{x_1:\omega_1,\dotsc,x_m:\omega_m} \\*
		\EnvSym\then \forall_{i \in [1,n]}{\tau_i \narreq \omega_i}
	}{
		\EnvSym\then \theta \narreq \rho
	}
\]
A tuple type $\theta$ is narrower than the record type $\rho$ if the first $n$
fields of $\theta$ have a type $\tau$ narrower than the first $n$ fields type
$\omega$ of $\rho$ (fields names are ignored).

\vspace{\baselineskip}

\[
	\infrule{
		\EnvSym\then n \geq 0 \semiq \forall_{i\in[1,n]}{\tau_i:\Type} \semiq \theta = \bigtimes_{i \in [1,n]}{\tau_i}\\*
		\EnvSym\then \xi \in \Id \semi m \geq n \semi \forall_{j\in[1,m]}{x_j \in \Id, \omega_j:\Type} \semi \sigma = (\xi, \tuple{x_1:\omega_1,\dotsc,x_m:\omega_m}) \\*
		\EnvSym\then \forall_{i \in [1,n]}{\tau_i \narreq \omega_i}
	}{
		\EnvSym\then \theta \narreq \sigma
	}
\]
A tuple type $\theta$ is narrower than the struct type $\sigma$ if the first
$n$ fields of $\theta$ have a type $\tau$ narrower than the first $n$ fields
type $\omega$ of $\sigma$ (struct name and fields names are ignored).

\vspace{\baselineskip}

\clearpage

\paragraph{Record}

\[
	\infrule{
		\EnvSym\then n \geq 0 \semiq \forall_{i\in[1,n]}{x_i \in \mathcal{N} \subseteq \mathcal{M}, \tau_i:\Type} \semiq \rho_1 = \tuple{x_1:\tau_1,\dotsc,x_n:\tau_n} \\*
		\EnvSym\then m \geq n \semiq \forall_{j\in[1,m]}{y_j \in \mathcal{M} \subseteq \Id, \omega_j:\Type} \semiq \rho_2 = \tuple{y_1:\omega_1,\dotsc,y_m:\omega_m} \\*
		\EnvSym\then \forall_{x':\tau'\in\mathcal{N}}{\exists_{y':\omega'\in\mathcal{M}}{x'=y' \land \tau' \narreq \omega'}}
	}{
		\EnvSym\then \rho_1 \narreq \rho_2
	}
\]
A record type $\rho_1$ is narrower than a record type $\rho_2$ if, for all
fields $x'$ of type $\tau'$ in $\rho_1$, there is a field with the same name
$y'$ of type $\omega'$ in $\rho_2$, where $\tau'$ is narrower than $\omega'$.

\vspace{\baselineskip}

\[
	\infrule{
		\EnvSym\then n \geq 0 \semiq \forall_{i\in[1,n]}{x \in \mathcal{N} \subseteq \mathcal{M}, \tau_i:\Type} \semiq \rho = \tuple{x_1:\tau_1,\dotsc,x_n:\tau_n} \\*
		\EnvSym\then \xi \in \Id \semi m \geq n \semi \forall_{j\in[1,m]}{y_j \in \Id, \omega_j:\Type} \semiq \sigma = (\xi, \tuple{y_1:\omega_1,\dotsc,y_m:\omega_m}) \\*
		\EnvSym\then \forall_{x':\tau'\in\mathcal{N}}{\exists_{y':\omega'\in\mathcal{M}}{x'=y' \land \tau' \narreq \omega'}}
	}{
		\EnvSym\then \rho \narreq \sigma
	}
\]
A record type $\rho$ is narrower than a struct type $\sigma$ if, for all fields
$x'$ of type $\tau'$ in $\rho$, there is a field with the same name $y'$ of
type $\omega'$ in $\sigma$, where $\tau'$ is narrower than $\omega'$ (struct
name is ignored).

\paragraph{Struct}

\[
	\infrule{
		\EnvSym\then \xi_1 \in \Id \semiq n \geq 0 \semiq \forall_{i\in[1,n]}{x_i \in \Id, \tau_i:\Type} \\*
		\EnvSym\then \sigma_1 = (\xi_1, \tuple{x_1:\tau_1,\dotsc,x_n:\tau_n}) \\*
		\EnvSym\then \xi_2 \in \Id \semiq m \geq n \semiq \forall_{j\in[1,m]}{y_j \in \Id, \omega_j:\Type} \\*
		\EnvSym\then \sigma_2 = (\xi_2, \tuple{y_1:\omega_1,\dotsc,y_m:\omega_m}) \\*
		\EnvSym\then \xi_1 = \xi_2
	}{
		\EnvSym\then \sigma_1 \narreq \sigma_2
	}
\]
A struct type $\sigma_1$ is narrower than a struct type $\sigma_2$ if and only
if they have the same name (fields are ignored since they're equal by
construction).

\todo{questa regola è ridondante?}
\vspace{\baselineskip}

\clearpage

\subsubsection{Function}

\[
	\infrule{
		\EnvSym\then n \geq 0 \semiq \forall_{i\in[1,n]}{x_i\in\mathcal{N}\subseteq\mathcal{M},\tau_i:\Type} \semiq \tau_{n+1}:\Type  \\*
		\EnvSym\then \phi_1 = (x_1:\tau_1,\dotsc,x_n:\tau_n)\to\tau_{n+1} \\*
		\EnvSym\then m \geq n \semiq \forall_{j\in[1,m]}{y_j\in\mathcal{M}\subseteq\Id,\omega_i:\Type} \semiq \omega_{m+1}:\Type \\*
		\EnvSym\then \phi_2 = (y_1:\omega_1,\dotsc,y_m:\omega_m)\to\omega_{m+1} \\*
		\EnvSym\then \tau_{n+1} \narreq \omega_{m+1} \\*
		\EnvSym\then \forall_{x':\tau'}{\exists_{y':\omega'}{x'=y' \land \tau' \narreq \omega'}}
	}{
		\EnvSym\then \phi_1 \narreq \phi_2
	}
\]
A function type $\phi_1$ is narrower than a function type $\phi_2$ if the
return type $\tau_{n+1}$ is narrower than the other return type $\omega_{m+1}$,
and if, for all parameters $x'$ of type $\tau'$ in $\phi_1$, there is a
parameter with the same name $y'$ of type $\omega'$ in $\phi_2$, where $\tau'$
is narrower than $\omega'$.

\todo{
	è gia un casino cosi :D, ma includere anche opzionalità dei parametri\\
	Eventualmente estrarre una regola per il singolo parametro
}


\commentinline{
	TODO: come menzionato nel README, il widening tra funzioni potrebbe essere
	errato se si considerano i domini delle funzioni rispetto al loro corpo.

	Per funzioni senza parametri è ok: \code{::() => Int} può essere allargata
	a \code{::() => Real}.

	Ma se una funziona richiede dei parametri \code{::(i Int) => Int},
	significa che il suo corpo sa come lavorare con \code{i Int}; nel momento
	in cui si allarga a \code{::(i Real) => Real}, il corpo potrebbe operare su
	un tipo inaspettato
}

\clearpage

\section{Expressions}

\subsection{Literals}

\subsubsection{Bool}

\[
	\EnvSym\then \code{true}:\Bool \qquad \EnvSym \then \code{false}:\Bool
\]

\subsubsection{Int}

\[
	\infrule{
		\EnvSym\then x \in \mathbb{Z} \semiq x \in [-2^{31}, 2^{31} - 1]
	}{
		\EnvSym\then \code{$x$} : \Int \\*
		\land \EnvSym\then \code{$x$i} : \Int
	}
\]

% \subparagraph{Examples} \verb|0, 0i, 69, -420, 32i, -99i|

\subsubsection{Natural}

\[
	\infrule{
		\EnvSym\then x \in \mathbb{N} \semiq x \in [0, 2^{31}]
	}{
		\EnvSym\then \code{$x$n} : \Natural
	}
\]

% \subparagraph{Examples} \verb|0n, 69n, 420n|

\subsubsection{Real}

\[
	\infrule{
		\EnvSym\then x \in \mathbb{R}
	}{
		\EnvSym\then \code{$x$} : \Real \\*
		\land \EnvSym\then \code{$x$r} : \Real
	}
\]

% \subparagraph{Examples} \verb|0r, 0.0, 6.9, -4.20|

\subsubsection{Text}

\[
	\infrule{
		\EnvSym\then x \in \text{Unicode} \setminus \set{"\backslash n", "\backslash r"} \semiq n \geq 0
	}{
		\EnvSym\then \code{"$x_1 \dotsc x_n$"} : \Text \\*
		\land \EnvSym\then \code{'$x_1 \dotsc x_n$'} : \Text
	}
\]

\todo{specificare il set $Unicode$ con i caratteri permessi da Java}

\[
	\infrule{
		\EnvSym\then x \in \text{Unicode} \semiq n \geq 0
	}{
		\EnvSym\then \code{"\{$x_1 \dotsc x_n$\}"} : \Text \\*
		\land \EnvSym\then \code{'\{$x_1 \dotsc x_n$\}'} : \Text
	}
\]

% \subparagraph{Examples} \verb|"hello", 'world', "{hello}", '{world}'|

\clearpage

\subsection{Collection types}

\subsubsection{List}

\[
	\EnvSym\then \code{[]}:\exists\Void.\List
\]

\[
	\infrule{
		\EnvSym\then n \geq 1 \semiq \tau:\Type \semiq \forall_{i\in[1,n]}{\expr_i\in\Expr,\expr_i:\tau}
	}{
		\EnvSym\then \code{[$\expr_1$, $\dotsc$, $\expr_n$]}:\exists\tau.\List
	}
\]

\[
	\infrule{
		\EnvSym\then n \geq 1 \semiq \forall_{i\in[1,n]}{\tau_i:\Type,\expr_i\in\Expr,\expr_i:\tau_i} \semiq \omega = \biglwb_{i=1}^{n}{\tau_i}
	}{
		\EnvSym\then \code{[$\expr_1$, $\dotsc$, $\expr_n$]}:\exists\omega.\List
	}
\]

List of $\omega$ (widest among $\tau$)

\subsubsection{Map}

\[
	\EnvSym\then \code{[=>]}:\exists\Void.\exists\Void.\Map
\]

\[
	\infrule{
		\EnvSym\then \kappa:\Type \semiq \forall_{i\in[1,n]}{k_i\in\Expr,k_i:\kappa} \semiq n \geq 1 \\*
		\EnvSym\then \nu:\Type \semiq \forall_{i\in[1,n]}{v_i\in\Expr,v_i:\nu} \hfill
	}{
		\EnvSym\then \code{[$k_1$=>$v_1$, $\dotsc$, $k_n$=>$v_n$]}:\exists\kappa.\exists\nu.\Map
	}
\]

\[
	\infrule{
		\EnvSym\then \forall_{i \in [1,n]}{\kappa_i:\Type, k_i\in\Expr, k_i:\kappa_i} \semiq n \geq 1 \\*
		\EnvSym\then \forall_{i \in [1,n]}{\nu_i:\Type, v_i\in\Expr, v_i:\nu_i} \hfill\\*
		\EnvSym\then \omega_\kappa = \biglwb_{i=1}^{n}{\kappa_i} \semiq \omega_\nu = \biglwb_{i=1}^{n}{\nu_i}
	}{
		\EnvSym\then \code{[$k_1$=>$v_1$, $\dotsc$, $k_n$=>$v_n$]}:\exists\omega_\kappa.\exists\omega_\nu.\Map
	}
\]

Map of $\omega_\kappa$ keys (widest  among $\kappa$) and $\omega_\nu$ (widest
among $\nu$).

\clearpage

\subsection{Product types}

\subsubsection{Tuple}

\[
	\infrule{
		\EnvSym\then n \geq 0 \semiq \forall_{i\in[1,n]}{\tau_i:\Type, \expr_i\in\Expr, \expr_i:\tau_i}
	}{
		\EnvSym\then \code{<$\expr_1$, $\dotsc$, $\expr_n$>}:\tau_1 \times \dotsc \times \tau_n
	}
\]

\subsubsection{Record}

\[
	\infrule{
		\EnvSym\then n \geq 1 \semiq \forall_{i\in[1,n]}{x_i\in\Id, \tau_i:\Type, \expr_i\in\Expr, \expr_i:\tau_i}
	}{
		\EnvSym\then \code{<$x_1$:$\expr_1$, $\dotsc$, $x_n$:$\expr_n$>}:\tuple{x_1:\tau_1,\dotsc,x_2:\tau_n}
	}
\]

\subsubsection{Struct}

\[
	\infrule{
		\EnvSym\then \xi \in \Id \semiq n \geq 0 \semiq \forall_{i\in[1,n]}{x_i \in \Id, \tau_i:\Type} \semiq \sigma = (\xi, \tuple{x_1:\tau_1,\dotsc,x_n:\tau_n}) \\*
	}{
		\EnvSym\then \code{$\xi$.<>} : \sigma \\*
	}
\]
This expression initializes a struct $\xi$ of type $\sigma$.

\[
	\infrule{
		\EnvSym\then \xi \in \Id \semiq n \geq 0 \semiq \forall_{i\in[1,n]}{x_i \in \Id, \tau_i:\Type} \semiq \sigma = (\xi, \tuple{x_1:\tau_1,\dotsc,x_n:\tau_n}) \\*
		\EnvSym\then \forall_{i\in[1,n]}{\eta_i:\Type,\eta_i\narreq\tau_i,\expr_i:\eta_i} \\*
	}{
		\EnvSym\then \code{$\xi$.<$\expr_1$, $\dotsc$, $\expr_n$>} : \sigma \\*
	}
\]
This expression initializes a struct $\xi$ of type $\sigma$ by using a
tuple-like notation, where each field gets initialized with the expression
$\expr_i$ of type $\eta_i$ (narrower than $\tau_i$).

\[
	\infrule{
		\EnvSym\then \xi \in \Id \semi n \geq 0 \semi \forall_{i\in[1,n]}{x_i\in\mathcal{N}\subseteq\Id, \tau_i:\Type} \semi \sigma = (\xi, \tuple{x_1:\tau_1,\dotsc,x_n:\tau_n}) \\*
		\EnvSym\then m \leq n \semi \forall_{j\in[1,m]}{y_j\in\mathcal{M}\subseteq\mathcal{N}, \eta_j:\Type, \expr_j:\eta_j} \\*
		\EnvSym\then \forall_{j\in[1,m]}{\exists_{i\in[1,n]}{y_j=x_i \land \eta_j\narreq\tau_i}}
		% \EnvSym\then \forall_{\mathcal{M}\ni y':\eta'}{\exists_{\mathcal{N}\ni x':\tau'}{y'=x' \land \eta'\narr\tau'}}
	}{
		\EnvSym\then \code{$\xi$.<$y_1$:$\expr_1$, $\dotsc$, $y_n$:$\expr_n$>} : \sigma \\*
	}
\]
This expression initializes a struct $\xi$ of type $\sigma$ by using a
record-like notation, where each specified field $y_j$ gets initialized with the
expression $\expr_j$ of type $\eta_j$ (narrower than $\tau_i$).

\clearpage

\subsection{Parenthesized expression}

\[
	\infrule{
		\EnvSym\then \tau:\Type \semiq \expr\in\Expr \semiq \expr:\tau
	}{
		\EnvSym\then \code{\{$\expr$\}} : \tau
	}
\]

\subsection{Invocation}

\paragraph{Positional}

\[
	\infrule{
		\EnvSym\then n \geq 0 \semiq \forall_{i\in[1,n]}{\tau_i:\Type} \semiq \tau_{n+1}:\Type \\*
		\EnvSym\then \phi = (\tau_1,\dotsc,\tau_n)\to\tau_{n+1} \semiq  f : \phi \\*
		\EnvSym\then \forall_{i\in[1,n]}{\eta_i:\Type,\eta_i\narreq\tau_i,\expr_i:\eta_i}
	}{
		\EnvSym\then \code{$f$($\expr_1$,$\dotsc$,$\expr_n$)} : \tau_{n+1}
	}
\]

\paragraph{Nominal}

\[
	\infrule{
		\EnvSym\then n \geq 0 \semiq \forall_{i\in[1,n]}{x\in\mathcal{N}\subseteq\Id,\tau_i:\Type} \semiq \tau_{n+1}:\Type \\*
		\EnvSym\then \phi = (x_1:\tau_1,\dotsc,x_n:\tau_n)\to\tau_{n+1} \semiq  f : \phi \\*
		\EnvSym\then \forall_{j\in[1,n]}{y\in\mathcal{N},\eta_j:\Type,\expr_j:\eta_j} \\*
		\EnvSym\then \forall_{j\in[1,n]}{\exists_{i\in[1,n]}{y_j=x_i,\eta_j\narreq\tau_i}}
	}{
		\EnvSym\then \code{$f$($y_1$:$\expr_1$,$\dotsc$,$y_n$:$\expr_n$)} : \tau_{n+1}
	}
\]

\todo{completare}

\subsection{Reference}

\todo{chiarire notazione e muovere in sezione operatori}

\[
	\infrule{
		\EnvSym\then x \in \Id \\*
		\EnvSym\then n \geq 0 \semiq \forall_{i\in[1,n]}{x_i\in\Id,\tau_i:\Type} \semiq \tau_{n+1}\Type \\*
		\EnvSym\then \phi = (x_1:\tau_1,\dotsc,x_n:\tau_n)\to\tau_{n+1} \\*
	}{
		\EnvSym,\code{$x$ $\phi$}\then \code{::$x$} : \phi
	}
\]

\clearpage

\subsection{Range}

\[
	\infrule{
		\EnvSym\then \tau_1 \in \Number \semiq \expr_1\in\Expr \semiq \expr_1:\tau_1\\*
		\EnvSym\then \tau_2 \in \Number \semiq \expr_2\in\Expr \semiq \expr_2:\tau_2 \\*
		\EnvSym\then \omega = \tau_1 \lwb \tau_2
	}{
		\EnvSym\then \code{$\expr_1$..$\expr_2$} : \exists\omega.\Range
	}
\]
Range from $\expr_1$ (begin) to $\expr_2$ (end).

\[
	\infrule{
		\EnvSym\then \tau_1 \in \Number \semiq \expr_1\in\Expr \semiq \expr_1:\tau_1\\*
		\EnvSym\then \tau_2 \in \Number \semiq \expr_2\in\Expr \semiq \expr_2:\tau_2 \\*
		\EnvSym\then \tau_3 \in \Number \semiq \expr_3\in\Expr \semiq \expr_3:\tau_3 \\*
		\EnvSym\then \omega = \tau_1 \lwb \tau_2 \lwb \tau_3
	}{
		\EnvSym\then \code{$\expr_1$..$\expr_2$..$\expr_3$} : \exists\omega.\Range
	}
\]
Range from $\expr_1$ (begin) to $\expr_2$ (end) with a step $\expr_3$.

\clearpage

\subsection{General rules}

\[
	\infrule{
		\EnvSym\then \tau_1:\Type \semiq \expr_1:\tau \\*
		\EnvSym\then \tau_2:\Type \semiq \tau_1 \narreq \tau_2 \semiq \expr_2:\tau_2
	}{
		\EnvSym\then \code{$\expr_1$ = $\expr_2$} : \Bool
	}
\]

\[
	\infrule{
		\EnvSym\then \tau_1:\Type \semiq \expr_1:\tau_1 \\*
		\EnvSym\then \tau_2:\Type \semiq \expr_2:\tau_2 \\*
		\EnvSym\then \code{$\expr_1$ = $\expr_2$} : \Bool
	}{
		\EnvSym\then \code{$\expr_1$ != $\expr_2$} : \Bool
	}
\]
If the "equal" operator is defined, the "not equal" operator is also defined.

\vspace{\baselineskip}

\[
	\infrule{
		\EnvSym\then \tau_1:\Type \semiq \expr_1:\tau_1 \\*
		\EnvSym\then \tau_2:\Type \semiq \expr_2:\tau_2 \\*
		\EnvSym\then \code{$\expr_1$ =\~{} $\expr_2$}: \Bool
	}{
		\EnvSym\then \code{$\expr_2$ \~{}= $\expr_1$}: \Bool
	}
\]
If the "contains" operator is defined, the "contained" operator is also
defined.

\vspace{\baselineskip}

\[
	\Number = \set{\Natural, \Int, \Real}
\]

\[
	\infrule{
		\EnvSym\then \expr:\Number \semiq \tau \in \Number
	}{
		\EnvSym\then \code{$\expr$ as $\tau$}: \tau
	}
\]

\todo[impl]{-69 as Natural fallisce (out of bounds). Dovrebbe?}

\[
	\infrule{
		\EnvSym\then \tau \in \Number \semiq \expr:\tau
	}{
		\EnvSym\then \code{+$\expr$}: \tau
	}
\]

\[
	\infrule{
		\EnvSym\then \tau \in \Number\setminus\set{\Natural} \semiq \expr:\tau
	}{
		\EnvSym\then \code{-$\expr$}: \tau
	}
\]

\[
	\infrule{
		\EnvSym\then \tau_1,\tau_2 \in \Number \semiq \expr_1:\tau_1 \semiq \expr_2:\tau_2\\*
		\bop = \set{\code{=}, \code{!=}, \code{<}, \code{<=}, \code{>=}, \code{>=}, \code{>}}
	}{
		\EnvSym\then \code{$\expr_1$ $\bop$ $\expr_2$}: \Bool
	}
\]

\clearpage

\subsection{Primitive types}

\subsubsection{Void}

\todo{?}

\subsubsection{Bool}

\[
	\infrule{
		\EnvSym\then \expr:\Bool
	}{
		\EnvSym\then \code{!$\expr$}:\Bool
	}
\]

\[
	\infrule{
		\EnvSym\then \expr_1:\Bool \semiq \expr_2:\Bool \\*
		\bop = \set{\code{and}, \code{\&}, \code{or}, \code{|}}
	}{
		\EnvSym\then \code{$\expr_1$ $\bop$ $\expr_2$}:\Bool
	}
\]

\subsubsection{Natural}

\paragraph{Plus}

\[
	\infrule{
		\EnvSym\then \expr_1:\Natural \semiq \expr_2:\Natural
	}{
		\EnvSym\then \code{$\expr_1$ + $\expr_2$} : \Natural
	}
\]

\[
	\infrule{
		\EnvSym\then \expr_1:\Natural \semiq \expr_2:\Int
	}{
		\EnvSym\then \code{$\expr_1$ + $\expr_2$} : \Int
	}
\]

\[
	\infrule{
		\EnvSym\then \expr_1:\Natural \semiq \expr_2:\Real
	}{
		\EnvSym\then \code{$\expr_1$ + $\expr_2$} : \Real
	}
\]

\paragraph{Minus}

\[
	\infrule{
		\EnvSym\then \expr_1:\Natural \semiq \expr_2:\Natural
	}{
		\EnvSym\then \code{$\expr_1$ - $\expr_2$} : \Int
	}
\]

\todo{specificare se $l-r\geq0$}

\[
	\infrule{
		\EnvSym\then \expr_1:\Natural \semiq \expr_2:\Int
	}{
		\EnvSym\then \code{$\expr_1$ - $\expr_2$} : \Int
	}
\]

\[
	\infrule{
		\EnvSym\then \expr_1:\Natural \semiq \expr_2:\Real
	}{
		\EnvSym\then \code{$\expr_1$ - $\expr_2$} : \Real
	}
\]

\paragraph{Multiplication}

\[
	\infrule{
		\EnvSym\then \expr_1:\Natural \semiq \expr_2:\Natural
	}{
		\EnvSym\then \code{$\expr_1$ * $\expr_2$} : \Natural
	}
\]


\[
	\infrule{
		\EnvSym\then \expr_1:\Natural \semiq \expr_2:\Int
	}{
		\EnvSym\then \code{$\expr_1$ * $\expr_2$} : \Int
	}
\]

\[
	\infrule{
		\EnvSym\then \expr_1:\Natural \semiq \expr_2:\Real
	}{
		\EnvSym\then \code{$\expr_1$ * $\expr_2$} : \Real
	}
\]

\paragraph{Division}

\[
	\infrule{
		\EnvSym\then \expr_1:\Natural \semiq \expr_2:\Natural
	}{
		\EnvSym\then \code{$\expr_1$ / $\expr_2$} : \Real
	}
\]

\[
	\infrule{
		\EnvSym\then \expr_1:\Natural \semiq \expr_2:\Int
	}{
		\EnvSym\then \code{$\expr_1$ / $\expr_2$} : \Real
	}
\]

\[
	\infrule{
		\EnvSym\then \expr_1:\Natural \semiq \expr_2:\Real
	}{
		\EnvSym\then \code{$\expr_1$ / $\expr_2$} : \Real
	}
\]

\paragraph{Modulo}

\[
	\infrule{
		\EnvSym\then \expr_1:\Natural \semiq \expr_2:\Natural
	}{
		\EnvSym\then \code{$\expr_1$ \% $\expr_2$} : \Natural
	}
\]

\[
	\infrule{
		\EnvSym\then \expr_1:\Natural \semiq \expr_2:\Int
	}{
		\EnvSym\then \code{$\expr_1$ \% $\expr_2$} : \Int
	}
\]

\[
	\infrule{
		\EnvSym\then \expr_1:\Natural \semiq \expr_2:\Real
	}{
		\EnvSym\then \code{$\expr_1$ \% $\expr_2$} : \Real
	}
\]

\paragraph{Power}

\[
	\infrule{
		\EnvSym\then \expr_1:\Natural \semiq \expr_2:\Natural
	}{
		\EnvSym\then \code{$\expr_1$ \^{} $\expr_2$} : \Natural
	}
\]

\[
	\infrule{
		\EnvSym\then \expr_1:\Natural \semiq \expr_2:\Int
	}{
		\EnvSym\then \code{$\expr_1$ \^{} $\expr_2$} : \Int
	}
\]

\[
	\infrule{
		\EnvSym\then \expr_1:\Natural \semiq \expr_2:\Real
	}{
		\EnvSym\then \code{$\expr_1$ \^{} $\expr_2$} : \Real
	}
\]

\subsubsection{Int}

\paragraph{Bitwise not}

\[
	\infrule{
		\EnvSym\then \expr:\Int
	}{
		\EnvSym\then \code{!!$\expr$} : \Int
	}
\]

\paragraph{Bitwise and}

\[
	\infrule{
		\EnvSym\then \expr_1:\Int \semiq \expr_2:\Natural
	}{
		\EnvSym\then \code{$\expr_1$ \&\& $\expr_2$} : \Int
	}
\]

\[
	\infrule{
		\EnvSym\then \expr_1:\Int \semiq \expr_2:\Int
	}{
		\EnvSym\then \code{$\expr_1$ \&\& $\expr_2$} : \Int
	}
\]

\paragraph{Bitwise or}

\[
	\infrule{
		\EnvSym\then \expr_1:\Int \semiq \expr_2:\Natural
	}{
		\EnvSym\then \code{$\expr_1$ || $\expr_2$} : \Int
	}
\]

\[
	\infrule{
		\EnvSym\then \expr_1:\Int \semiq \expr_2:\Int
	}{
		\EnvSym\then \code{$\expr_1$ || $\expr_2$} : \Int
	}
\]

\paragraph{Bitwise xor}

\[
	\infrule{
		\EnvSym\then \expr_1:\Int \semiq \expr_2:\Natural
	}{
		\EnvSym\then \code{$\expr_1$ ||| $\expr_2$} : \Int
	}
\]

\[
	\infrule{
		\EnvSym\then \expr_1:\Int \semiq \expr_2:\Int
	}{
		\EnvSym\then \code{$\expr_1$ ||| $\expr_2$} : \Int
	}
\]

\paragraph{Bitwise left shift}

\[
	\infrule{
		\EnvSym\then \expr_1:\Int \semiq \expr_2:\Natural
	}{
		\EnvSym\then \code{$\expr_1$ \{<\} $\expr_2$} : \Int
	}
\]

\[
	\infrule{
		\EnvSym\then \expr_1:\Int \semiq \expr_2:\Int
	}{
		\EnvSym\then \code{$\expr_1$ \{<\} $\expr_2$} : \Int
	}
\]

\paragraph{Bitwise right shift}

\[
	\infrule{
		\EnvSym\then \expr_1:\Int \semiq \expr_2:\Natural
	}{
		\EnvSym\then \code{$\expr_1$ \{>\} $\expr_2$} : \Int
	}
\]

\[
	\infrule{
		\EnvSym\then \expr_1:\Int \semiq \expr_2:\Int
	}{
		\EnvSym\then \code{$\expr_1$ \{>\} $\expr_2$} : \Int
	}
\]

\paragraph{Bitwise right unsigned shift}

\[
	\infrule{
		\EnvSym\then \expr_1:\Int \semiq \expr_2:\Natural
	}{
		\EnvSym\then \code{$\expr_1$ \{>>\} $\expr_2$} : \Int
	}
\]

\[
	\infrule{
		\EnvSym\then \expr_1:\Int \semiq \expr_2:\Int
	}{
		\EnvSym\then \code{$\expr_1$ \{>>\} $\expr_2$} : \Int
	}
\]

\paragraph{Plus}

\[
	\infrule{
		\EnvSym\then \expr_1:\Int \semiq \expr_2:\Natural
	}{
		\EnvSym\then \code{$\expr_1$ + $\expr_2$} : \Int
	}
\]

\[
	\infrule{
		\EnvSym\then \expr_1:\Int \semiq \expr_2:\Int
	}{
		\EnvSym\then \code{$\expr_1$ + $\expr_2$} : \Int
	}
\]

\[
	\infrule{
		\EnvSym\then \expr_1:\Int \semiq \expr_2:\Real
	}{
		\EnvSym\then \code{$\expr_1$ + $\expr_2$} : \Real
	}
\]

\paragraph{Minus}

\[
	\infrule{
		\EnvSym\then \expr_1:\Int \semiq \expr_2:\Natural
	}{
		\EnvSym\then \code{$\expr_1$ - $\expr_2$} : \Int
	}
\]

\[
	\infrule{
		\EnvSym\then \expr_1:\Int \semiq \expr_2:\Int
	}{
		\EnvSym\then \code{$\expr_1$ - $\expr_2$} : \Int
	}
\]

\[
	\infrule{
		\EnvSym\then \expr_1:\Int \semiq \expr_2:\Real
	}{
		\EnvSym\then \code{$\expr_1$ - $\expr_2$} : \Real
	}
\]

\paragraph{Multiplication}

\[
	\infrule{
		\EnvSym\then \expr_1:\Int \semiq \expr_2:\Natural
	}{
		\EnvSym\then \code{$\expr_1$ * $\expr_2$} : \Int
	}
\]

\[
	\infrule{
		\EnvSym\then \expr_1:\Int \semiq \expr_2:\Int
	}{
		\EnvSym\then \code{$\expr_1$ * $\expr_2$} : \Int
	}
\]

\[
	\infrule{
		\EnvSym\then \expr_1:\Int \semiq \expr_2:\Real
	}{
		\EnvSym\then \code{$\expr_1$ * $\expr_2$} : \Real
	}
\]

\paragraph{Division}

\[
	\infrule{
		\EnvSym\then \expr_1:\Int \semiq \expr_2:\Natural
	}{
		\EnvSym\then \code{$\expr_1$ / $\expr_2$} : \Real
	}
\]

\[
	\infrule{
		\EnvSym\then \expr_1:\Int \semiq \expr_2:\Int
	}{
		\EnvSym\then \code{$\expr_1$ / $\expr_2$} : \Real
	}
\]

\[
	\infrule{
		\EnvSym\then \expr_1:\Int \semiq \expr_2:\Real
	}{
		\EnvSym\then \code{$\expr_1$ / $\expr_2$} : \Real
	}
\]

\paragraph{Modulo}

\[
	\infrule{
		\EnvSym\then \expr_1:\Int \semiq \expr_2:\Natural
	}{
		\EnvSym\then \code{$\expr_1$ \% $\expr_2$} : \Int
	}
\]

\[
	\infrule{
		\EnvSym\then \expr_1:\Int \semiq \expr_2:\Int
	}{
		\EnvSym\then \code{$\expr_1$ \% $\expr_2$} : \Int
	}
\]

\[
	\infrule{
		\EnvSym\then \expr_1:\Int \semiq \expr_2:\Real
	}{
		\EnvSym\then \code{$\expr_1$ \% $\expr_2$} : \Real
	}
\]

\paragraph{Power}

\[
	\infrule{
		\EnvSym\then \expr_1:\Int \semiq \expr_2:\Natural
	}{
		\EnvSym\then \code{$\expr_1$ \^{} $\expr_2$} : \Real
	}
\]

\[
	\infrule{
		\EnvSym\then \expr_1:\Int \semiq \expr_2:\Int
	}{
		\EnvSym\then \code{$\expr_1$ \^{} $\expr_2$} : \Real
	}
\]

\[
	\infrule{
		\EnvSym\then \expr_1:\Int \semiq \expr_2:\Real
	}{
		\EnvSym\then \code{$\expr_1$ \^{} $\expr_2$} : \Real
	}
\]

\subsubsection{Real}

\paragraph{Plus}

\[
	\infrule{
		\EnvSym\then \expr_1:\Int \semiq \tau \in \Number \semiq \expr_2:\tau
	}{
		\EnvSym\then \code{$\expr_1$ + $\expr_2$} : \Real
	}
\]

\paragraph{Minus}

\[
	\infrule{
		\EnvSym\then \expr_1:\Int \semiq \tau \in \Number \semiq \expr_2:\tau
	}{
		\EnvSym\then \code{$\expr_1$ - $\expr_2$} : \Real
	}
\]

\paragraph{Multiplication}

\[
	\infrule{
		\EnvSym\then \expr_1:\Int \semiq \tau \in \Number \semiq \expr_2:\tau
	}{
		\EnvSym\then \code{$\expr_1$ * $\expr_2$} : \Real
	}
\]

\paragraph{Division}

\[
	\infrule{
		\EnvSym\then \expr_1:\Int \semiq \tau \in \Number \semiq \expr_2:\tau
	}{
		\EnvSym\then \code{$\expr_1$ / $\expr_2$} : \Real
	}
\]

\paragraph{Modulo}

\[
	\infrule{
		\EnvSym\then \expr_1:\Int \semiq \tau \in \Number \semiq \expr_2:\tau
	}{
		\EnvSym\then \code{$\expr_1$ \% $\expr_2$} : \Real
	}
\]

\paragraph{Power}

\[
	\infrule{
		\EnvSym\then \expr_1:\Int \semiq \tau \in \Number \semiq \expr_2:\tau
	}{
		\EnvSym\then \code{$\expr_1$ \^{} $\expr_2$} : \Real
	}
\]

\subsubsection{Text}

\paragraph{As}

\[
	\infrule{
		\EnvSym\then \expr_1:\Text \semiq \tau \in \Number \cup \set{\Text, \Bool}
	}{
		\EnvSym\then \code{$\expr_1$ as $\tau$}: \tau
	}
\]

\paragraph{Plus}

\[
	\infrule{
		\EnvSym\then \expr_1:\Text \semiq \tau:\Type \semiq \expr_2:\tau
	}{
		\EnvSym\then \code{$\expr_1$ + $\expr_2$}: \Text
	}
\]

\paragraph{Minus}

\[
	\infrule{
		\EnvSym\then \expr_1:\Text \semiq \expr_2:\Text
	}{
		\EnvSym\then \code{$\expr_1$ - $\expr_2$}: \Text
	}
\]

\paragraph{Power}

\[
	\infrule{
		\EnvSym\then \expr_1:\Text \semiq \tau \in \Number \semiq \expr_2:\tau
	}{
		\EnvSym\then \code{$\expr_1$ \^{} $\expr_2$}: \Text
	}
\]

\todo[impl]{non dovremmo permettere Real}

\paragraph{Division}

\[
	\infrule{
		\EnvSym\then \expr_1:\Text \semiq \expr_2:\Text
	}{
		\EnvSym\then \code{$\expr_1$ / $\expr_2$}: \exists\Text.\List
	}
\]

\paragraph{Equal}

\[
	\infrule{
		\EnvSym\then \expr_1:\Text \semiq \expr_2:\Text
	}{
		\EnvSym\then \code{$\expr_1$ = $\expr_2$}: \Bool
	}
\]

\paragraph{Less}

\[
	\infrule{
		\EnvSym\then \expr_1:\Text \semiq \expr_2:\Text
	}{
		\EnvSym\then \code{$\expr_1$ < $\expr_2$}: \Bool
	}
\]

\paragraph{Less equal}

\[
	\infrule{
		\EnvSym\then \expr_1:\Text \semiq \expr_2:\Text
	}{
		\EnvSym\then \code{$\expr_1$ <= $\expr_2$}: \Bool
	}
\]

\paragraph{Greater equal}

\[
	\infrule{
		\EnvSym\then \expr_1:\Text \semiq \expr_2:\Text
	}{
		\EnvSym\then \code{$\expr_1$ >= $\expr_2$}: \Bool
	}
\]

\paragraph{Greater}

\[
	\infrule{
		\EnvSym\then \expr_1:\Text \semiq \expr_2:\Text
	}{
		\EnvSym\then \code{$\expr_1$ > $\expr_2$}: \Bool
	}
\]

\paragraph{Contains}

\[
	\infrule{
		\EnvSym\then \expr_1:\Text \semiq \expr_2:\Text
	}{
		\EnvSym\then \code{$\expr_1$ =\~{} $\expr_2$}: \Bool
	}
\]

\paragraph{Dot access}

\[
	\infrule{
		\text{Index access}\\*
		\EnvSym\then \expr_1:\Text \semiq \expr_2:\Int
	}{
		\EnvSym\then \code{$\expr_1$.$\expr_2$}: \Text
	}
\]

\[
	\infrule{
		\text{"Slice"}\\*
		\EnvSym\then \expr_1:\Text \semiq \expr_2:\exists\Int.\Range
	}{
		\EnvSym\then \code{$\expr_1$.$\expr_2$}: \Text
	}
\]

\[
	\infrule{
		\text{Text length}\\*
		\EnvSym\then \expr_1:\Text
	}{
		\EnvSym\then \code{$\expr_1$.\_count}: \Int
	}
\]

\[
	\infrule{
		\text{Characters range}\\*
		\EnvSym\then \expr_1:\Text
	}{
		\EnvSym\then \code{$\expr_1$.\_range}: \exists\Int.\Range
	}
\]

\[
	\infrule{
		\text{Characters list}\\*
		\EnvSym\then \expr_1:\Text
	}{
		\EnvSym\then \code{$\expr_1$.\_chars}: \exists\Text.\List
	}
\]

\clearpage

\subsection{Collection types}

\subsubsection{List}

\paragraph{As}

\[
	\infrule{
		\EnvSym\then \tau_1 = \List[\nu_1 := \Type] \\*
		\EnvSym\then \tau_2 = \List[\nu_2 := \Type] \\*
		\EnvSym\then \expr:\tau_1 \semiq x:\nu_1 \\*
		\EnvSym\then \code{$x$ as $\nu_2$} : \nu_2
	}{
		\EnvSym\then \code{$\expr$ as $\tau_2$} : \tau_2
	}
\]
If a term of type $\nu_1$ can be casted to type $\nu_2$ $\implies$ A list
of $\nu_1$ can be casted to a list of $\nu_2$.

\paragraph{Plus}

\[
	\infrule{
		\EnvSym\then \tau_1=\List \semiq \expr_1:\tau_1 \\*
		\EnvSym\then \tau_2=\List \semiq \expr_2:\tau_2 \\*
		\EnvSym\then \tau_2 \narreq \tau_1
	}{
		\EnvSym\then \code{$\expr_1$ + $\expr_2$} : \tau_1
	}
\]
\todo{chiarire. Il tipo della concatenazione dovrebbe essere deciso da
lhs o dovrebbe essere il widest tra i due?}

\[
	\infrule{
		\EnvSym\then \tau=\List[\nu_1:=\Type] \semiq \expr_1:\tau \\*
		\EnvSym\then \nu_2:\Type \semiq \nu_2\narreq\nu_1 \semiq \expr_2:\nu_2
	}{
		\EnvSym\then \code{$\expr_1$ + $\expr_2$} : \tau
	}
\]

\paragraph{Minus}

\[
	\infrule{
		\EnvSym\then \tau_1=\List \semiq \expr_1:\tau_1 \\*
		\EnvSym\then \tau_2=\List \semiq \expr_2:\tau_2 \\*
		\EnvSym\then \tau_2 \narreq \tau_1
	}{
		\EnvSym\then \code{$\expr_1$ - $\expr_2$} : \tau_1
	}
\]
\todo{chiarire. Il tipo dell'intersezione dovrebbe essere deciso da
lhs o dovrebbe essere il widest tra i due?}

\[
	\infrule{
		\EnvSym\then \tau=\List[\nu_1:=\Type] \semiq \expr_1:\tau \\*
		\EnvSym\then \nu_2:\Type \semiq \nu_2\narreq\nu_1 \semiq \expr_2:\nu_2
	}{
		\EnvSym\then \code{$\expr_1$ - $\expr_2$} : \tau
	}
\]

\paragraph{Contains}

\[
	\infrule{
		\EnvSym\then \tau_1:\Type \semiq \expr_1:\exists\tau_1.\List \\*
		\EnvSym\then \tau_2:\Type \semiq \tau_2 \narreq \tau_1 \semiq \expr_2:\tau_2
	}{
		\EnvSym\then \code{$\expr_1$ =\~{} $\expr_2$} : \Bool
	}
\]
A list of $\tau_1$ contains an element of type $\tau_2$ (narrower than
$\tau_1$).

\[
	\infrule{
		\EnvSym\then \tau_1:\Type \semiq \expr_1:\exists\tau_1.\List\\*
		\EnvSym\then \tau_2:\Type \semiq \tau_2 \narreq \tau_1 \semiq \expr_2:\exists\tau_2.\List
	}{
		\EnvSym\then \code{$\expr_1$ =\~{} $\expr_2$} : \Bool
	}
\]
A list of $\tau_1$ contains the elements of a list of $\tau_2$ (narrower than
$\tau_1$).

\paragraph{Dot access}

\[
	\infrule{
		\text{Index access}\\*
		\EnvSym\then \tau:\Type
		\semiq \expr_1:\exists\tau.\List
		\semiq \expr_2:\Int
	}{
		\EnvSym\then \code{$\expr_1$.$\expr_2$}: \tau
	}
\]

\[
	\infrule{
		\text{"Slice"}\\*
		\EnvSym\then \tau : \Type
		\semiq \expr_1:\exists\tau.\List
		\semiq \expr_2:\exists\Int.\Range
	}{
		\EnvSym\then \code{$\expr_1$.$\expr_2$}: \tau
	}
\]

\[
	\infrule{
		\text{Count access}\\*
		\EnvSym\then \tau:\Type \semiq \expr:\exists\tau.\List
	}{
		\EnvSym\then \code{$\expr$.\_count}: \Int
	}
\]

\[
	\infrule{
		\text{Range access}\\*
		\EnvSym\then \expr:\Text
	}{
		\EnvSym\then \code{$\expr$.\_range}: \exists\Int.\Range
	}
\]

\subsubsection{Map}

\paragraph{As}

\[
	\infrule{
		\EnvSym\then \tau_1 = \Map[\kappa_1:=\Type,\nu_1:=\Type] \\*
		\EnvSym\then \tau_2 = \Map[\kappa_2:=\Type,\nu_2:=\Type] \\*
		\EnvSym\then \expr:\tau_1 \semiq \expr_k:\kappa_1 \semiq \expr_v:\nu_1 \\*
		\EnvSym\then \code{$\expr_k$ as $\kappa_2$} : \kappa_2 \land \code{$\expr_v$ as $\nu_2$} : \nu_2
	}{
		\EnvSym\then \code{$\expr$ as $\tau_2$} : \tau_2
	}
\]
If a term of type $\kappa_1$ and be casted to type $\kappa_2$ $\land$ A term of
type $\nu_1$ can be casted to type $\nu_2$ $\implies$ A map of type $\tau_1$
can be casted to the type $\tau_2$.

\paragraph{Plus}

\[
	\infrule{
		\EnvSym\then \tau_1=\Map \semiq \expr_1:\tau_1 \\*
		\EnvSym\then \tau_2=\Map \semiq \expr_2:\tau_2 \\*
		\EnvSym\then \tau_2 \narreq \tau_1
	}{
		\EnvSym\then \code{$\expr_1$ + $\expr_2$} : \tau_1
	}
\]
\todo{chiarire. Il tipo della concatenazione dovrebbe essere deciso da
lhs o dovrebbe essere il widest tra i due?}

\paragraph{Minus}

\[
	\infrule{
		\EnvSym\then \tau_1=\Map \semiq \expr_1:\tau_1 \\*
		\EnvSym\then \tau_2=\Map \semiq \expr_2:\tau_2 \\*
		\EnvSym\then \tau_2 \narreq \tau_1
	}{
		\EnvSym\then \code{$\expr_1$ - $\expr_2$} : \tau_1
	}
\]
\todo{chiarire. Il tipo dell'intersezione dovrebbe essere deciso da
lhs o dovrebbe essere il widest tra i due?}

\[
	\infrule{
		\EnvSym\then \tau=\Map[\kappa_1:=\Type] \semiq \expr_1:\tau \\*
		\EnvSym\then \kappa_2:\Type \semiq \kappa_2\narreq\kappa_1 \semiq \expr_2:\kappa_2
	}{
		\EnvSym\then \code{$\expr_1$ - $\expr_2$} : \tau
	}
\]


\paragraph{Contains}

\[
	\infrule{
		\EnvSym\then \tau = \Map[\kappa_1:=\Type,\nu:=\Type] \semiq \expr_1:\tau \\*
		\EnvSym\then \kappa_2:\Type \semiq \kappa_2 \narreq \kappa_1 \semiq \expr_2:\kappa_2
	}{
		\EnvSym\then \code{$\expr_1$ =\~{} $\expr_2$} : \Bool
	}
\]
A map $\expr_1$ of type $\tau$ contains a key of type $\kappa_2$ (narrower than
$\kappa_1$).

\[
	\infrule{
		\EnvSym\then \tau_1 = \Map[\kappa_1:=\Type,\nu_1:=\Type] \semiq \expr_1:\tau_1 \\*
		\EnvSym\then \tau_2 = \Map[\kappa_2:=\Type,\nu_2:=\Type] \semiq \expr_2:\tau_2 \\*
		\EnvSym\then \kappa_2 \narreq \kappa_1 \land \nu_2 \narreq \nu_1
	}{
		\EnvSym\then \code{$\expr_1$ =\~{} $\expr_2$} : \Bool
	}
\]
A map $\expr_1$ of type $\tau_1$ contais the key-value pairs of type $\kappa_2$
(narrower than $\kappa_1$) and $\nu_2$ (narrower than $\nu_1$) of the map
$\expr_2$ of type $\tau_2$.

\paragraph{Dot access}

\[
	\infrule{
		\text{Access by key}\\*
		\EnvSym\then \tau = \Map[\kappa_1:=\Type,\nu:=\Type] \semiq \expr_1:\tau \\*
		\EnvSym\then \kappa_2 = \Type \semiq \kappa_2 \narreq \kappa_1 \semiq \expr_2:\kappa_2
	}{
		\EnvSym\then \code{$\expr_1$.$\expr_2$}: \nu
	}
\]

\[
	\infrule{
		\text{Count access}\\*
		\EnvSym\then \tau = \Map[\kappa:=\Type,\nu:=\Type] \semiq \expr:\tau \\*
	}{
		\EnvSym\then \code{$\expr$.\_count}: \Int
	}
\]

\[
	\infrule{
		\text{Keys access}\\*
		\EnvSym\then \tau = \Map[\kappa:=\Type,\nu:=\Type] \semiq \expr:\tau \\*
	}{
		\EnvSym\then \code{$\expr$.\_keys}: \exists\kappa.\List
	}
\]

\[
	\infrule{
		\text{Values access}\\*
		\EnvSym\then \tau = \Map[\kappa:=\Type,\nu:=\Type] \semiq \expr:\tau \\*
	}{
		\EnvSym\then \code{$\expr$.\_values}: \exists\nu.\List
	}
\]

\subsection{Product types}

\[
	\Product = \set{\texttt{Tuple, Record, Struct}}
\]

\[
	\infrule{
		\EnvSym\then \tau_1 \in \Product \semiq \expr_1:\tau_1 \\*
		\EnvSym\then \tau_2 \in \Product \semiq \expr_2:\tau_2
	}{
		\EnvSym\then \code{$\expr_1$ = $\expr_2$}: \Bool
	}
\]

\[
	\infrule{
		\text{Fields count access} \\*
		\EnvSym\then \tau\in\Product \semiq \expr:\tau
	}{
		\EnvSym\then \code{$\expr$.\_count}: \Int
	}
\]

\[
	\infrule{
		\text{Fields range access} \\*
		\EnvSym\then \tau\in\Product \semiq \expr:\tau
	}{
		\EnvSym\then \code{$\expr$.\_range}: \exists\Int.\Range
	}
\]

\[
	\infrule{
		\text{Fields names access} \\*
		\EnvSym\then \tau\in\set{\Record,\Struct} \semiq \expr:\tau
	}{
		\EnvSym\then \code{$\expr$.\_fields}: \exists\Text.\List
	}
\]

\subsubsection{Tuple}

\paragraph{Dot access}

\[
	\infrule{
		\text{"Slice"}\\*
		\EnvSym\then n \geq 0 \semiq \forall_{i\in[1,n]}{\tau_i:\Type} \semiq \expr_1:\tau_1 \times \dotsb \times \tau_n \\*
		\EnvSym\then \expr_2:\exists\Int.\Range
	}{
		\EnvSym\then \code{$\expr_1$.$\expr_2$}: \bigtimes_{i \in \expr_2} \tau_i
	}
\]

Example: \code{<0,1,2,3,4,5> . 2..4 = <2 3 4>}

With step: \code{<0,1,2,3,4,5> . 0..5..2 = <0 2 4>}

Reverse: \code{<0,1,2,3,4,5> . 4..2 = <4 3 2>}

\[
	\infrule{
		\text{Access by index}\\*
		\EnvSym\then n \geq 0 \semiq \forall_{i\in[1,n]}{\tau_i:\Type} \semiq \expr_1:\tau_1 \times \dotsb \times \tau_n \\*
		\EnvSym\then \expr_2:\Int \semiq \expr_2 \in [1,n]
	}{
		\EnvSym\then \code{$\expr_1$.$\expr_2$}: \tau_{\expr_2}
	}
\]

\subsubsection{Record}

\paragraph{Dot access}

\[
	\infrule{
		\text{Access by index} \\*
		\EnvSym\then n \geq 1 \semiq \forall_{i\in[1,n]}{\tau_i:\Type, x_i\in\Id} \semiq \expr_1:\tuple{x_1:\tau_1,\dotsc,x_n:\tau_n}\\*
		\EnvSym\then \expr_2:\Int \semiq \expr_2 \in [1,n]
	}{
		\EnvSym\then \code{$\expr_1$.$\expr_2$}: \tau_{\expr_2}
	}
\]

\todo[impl]{permettere solo a Tuple in modo da rendere esplicito la perdita degli identificatori}

\[
	\infrule{
		\text{"Slice"}\\*
		\EnvSym\then n \geq 1 \semiq \forall_{i\in[1,n]}{\tau_i:\Type, x_i\in\Id} \semiq \rho=\tuple{x_1:\tau_1,\dotsc,x_n:\tau_n} \semiq \expr_1:\rho\\*
		\EnvSym\then \expr_2:\exists\Int.\Range \\*
		\EnvSym\then \forall_{i\in\expr_2}{\rho'=\tuple{x_j:\tau_j, \dotsc, x_j:\tau_j}}
	}{
		\EnvSym\then \code{$\expr_1$.$\expr_2$}: \rho'
	}
\]

\[
	\infrule{
		\text{Access by name}\\*
		\EnvSym\then n \geq 1 \semiq \forall_{i\in[1,n]}{\tau_i:\Type, x_i \in \Id} \semiq \expr:\tuple{x_1:\tau_1,\dotsc,x_n:\tau_n} \\*
		\EnvSym\then x' \in \Id \semiq \exists_{j\in[1,n]} x_j = x'
	}{
		\EnvSym\then \code{$\expr$.$x'$}: \tau_j
	}
\]

\subsubsection{Struct}

\todo{...}

\subsection{Function}

\todo{...}

\clearpage

\section{Statements}

\subsection{Variable declaration}

\[
	\infrule{
		\EnvSym\then x \in \Id \semiq \tau:\Type
	}{
		\EnvSym, \code{$x$ $\tau$} \then x:\tau
	}
\]

\[
	\infrule{
		\EnvSym\then x\in\Id \semiq \tau:\Type \\*
		\EnvSym\then \tau_\expr:\Type \semiq \tau_\expr\narreq\tau \semiq \expr\in\Expr \semiq \expr:\tau_\expr
	}{
		\EnvSym, \code{var $x$ $\tau$ <- $\expr$} \then x:\tau \qquad \text{(esplicit type)} \\*
		\land \EnvSym, \code{var $x$ <- $\expr$} \then x:\tau \qquad \text{(implicit type)}
	}
\]

\todo{specificare il tipo Var come da implementazione?}

\subsection{Constant declaration}

\[
	\infrule{
		\EnvSym\then x\in\Id \semiq \tau:\Type \\*
		\EnvSym\then \tau_\expr:\Type \semiq \tau_\expr\narreq\tau \semiq \expr\in\Expr \semiq \expr:\tau_\expr
	}{
		\EnvSym, \code{const $x$ $\tau$ <- $\expr$} \then x:\tau \qquad \text{(esplicit type)} \\*
		\land \EnvSym, \code{const $x$ <- $\expr$} \then x:\tau \qquad \text{(implicit type)}
	}
\]

\subsection{Alias declaration}

\[
	\infrule{
		\EnvSym\then x \in \Id \semiq \expr \in \Expr \semiq \tau:\Type \semiq \expr:\tau
	}{
		\EnvSym, \code{alias $e$ as $x$} \then x:\tau
	}
\]

\subsection{For loop}

\[
	\infrule{
		\EnvSym\then x \in \Id \semiq \tau:\Type \semiq \expr:\exists\tau.\List
	}{
		\EnvSym, \code{for $x$ in $\expr$ do} \then x:\tau
	}
\]

\[
	\infrule{
		\EnvSym\then k \in \Id \semiq v \in \Id \semiq \kappa:\Type \semiq \nu:\Type \semiq \expr:\exists\kappa.\exists\nu.\Map
	}{
		\EnvSym, \code{for $k$=>$v$ in $\expr$ do} \then k:\kappa \semiq v:\nu
	}
\]

\todo{deconstruction}

